\documentclass[11pt]{article}
\usepackage[a4paper,margin=1in]{geometry}
\usepackage{amsmath,amssymb,amsthm,mathtools}

\newcommand{\ones}{\mathbf{1}}
\newcommand{\diag}{\operatorname{diag}}
\newcommand{\norm}[1]{\left\lVert #1\right\rVert}
\newtheorem{proposition}{Proposition}

\title{Asymptotic Error Expansions for Forward/Backward Discrete Diffusion on Finite Markov Chains}
\author{}
\date{}

\begin{document}
\maketitle

\section{Setup}
Let $P\in\mathbb R^{n\times n}$ be a row-stochastic irreducible aperiodic Markov kernel, and define
\[
M\coloneqq P^\top.
\]
Forward histograms evolve by
\[
h^{t+1}=Mh^t,
\qquad t\in\mathbb N,
\]
with unique stationary histogram $h^\infty\in\Delta_n$ such that
\[
Mh^\infty=h^\infty.
\]

\section{Warm-up: expansion of $h^T-h^\infty$ for large $T$}
Assume $M$ is diagonalizable over $\mathbb R$ (or use Jordan form with the same conclusion up to polynomial factors). Write
\[
M = \Pi + \sum_{k=2}^n \lambda_k v_k w_k^\top,
\qquad
\Pi\coloneqq h^\infty \ones^\top,
\]
with $|\lambda_2|\ge\cdots\ge|\lambda_n|$ and $|\lambda_k|<1$ for $k\ge2$.
Then
\[
h^T-h^\infty
= M^T(h^0-h^\infty)
= \sum_{k=2}^n a_k\lambda_k^T v_k,
\qquad
 a_k\coloneqq w_k^\top(h^0-h^\infty).
\]
Hence the asymptotic expansion is
\[
h^T-h^\infty
= a_2\lambda_2^T v_2 + \mathcal O\big(|\lambda_3|^T\big).
\]
So the terminal forward bias is exponentially small in $T$, at rate $|\lambda_2|$.

\paragraph{Equivalent continuous-time notation.}
If $\lambda_k=e^{-\gamma_k}$, then
\[
h^T-h^\infty
= a_2e^{-\gamma_2 T}v_2+\mathcal O(e^{-\gamma_3T}).
\]

\section{Bayes reverse kernels and backward error expansion}
Define the Bayes reverse kernel (for each $t$)
\[
Q^t_{ij}=\frac{h_j^t P_{ji}}{h_i^{t+1}},
\qquad i,j\in\{1,\dots,n\},
\]
and set
\[
R_t\coloneqq (Q^t)^\top.
\]
Equivalently,
\[
R_t=\diag(h^t)\,P\,\diag(h^{t+1})^{-1}.
\]
Then exact backward marginal recursion is
\[
h^t = R_t h^{t+1},
\qquad t=T-1,\dots,0.
\]
Now define the perturbed backward recursion initialized at stationarity,
\[
g^T=h^\infty,
\qquad
g^t=R_t g^{t+1},\quad t=T-1,\dots,0.
\]
The error $e^t\coloneqq h^t-g^t$ satisfies
\[
e^t=R_t e^{t+1},
\qquad
 e^T=h^T-h^\infty,
\]
so
\[
e^0 = \Big(\prod_{t=0}^{T-1}R_t\Big)(h^T-h^\infty).
\]
Insert the warm-up expansion:
\[
e^0 = \Big(\prod_{t=0}^{T-1}R_t\Big)
\left(\sum_{k=2}^n a_k\lambda_k^T v_k\right).
\]
This is exact.

\subsection*{Explicit limit operator and fine asymptotics}
\begin{proposition}[Closed form of $R_\infty$]
Assume $h^\infty_i>0$ for all $i$. Then
\[
R_\infty
\coloneqq \lim_{t\to\infty}R_t
=\diag(h^\infty)\,P\,\diag(h^\infty)^{-1}.
\]
Hence $R_\infty$ is similar to $P$ and has the same spectrum.
\end{proposition}
\begin{proof}
From $R_t=\diag(h^t)P\diag(h^{t+1})^{-1}$ and $h^t\to h^\infty$.
\end{proof}

Write $\delta^t\coloneqq h^t-h^\infty$, so $\delta^t=\mathcal O(|\lambda_2|^t)$. With
\[
E_t\coloneqq \diag(\delta^t),\qquad D_\infty\coloneqq \diag(h^\infty),
\]
we have
\[
R_t=(D_\infty+E_t)P(D_\infty+E_{t+1})^{-1}
=R_\infty + B_t + \mathcal O(\norm{\delta^t}^2+\norm{\delta^{t+1}}^2),
\]
with first-order term
\[
B_t = E_t P D_\infty^{-1}-D_\infty P D_\infty^{-1}E_{t+1}D_\infty^{-1},
\qquad
\norm{B_t}=\mathcal O(|\lambda_2|^t).
\]

For the product, first-order Duhamel expansion gives
\[
\prod_{t=0}^{T-1}R_t
=R_\infty^T
+\sum_{s=0}^{T-1}R_\infty^{\,s}B_sR_\infty^{\,T-1-s}
+\mathcal O\!\Big(\sum_{s=0}^{T-1}|\lambda_2|^{2s}\Big).
\]
Therefore
\[
e^0
=R_\infty^T\delta^T
+\sum_{s=0}^{T-1}R_\infty^{\,s}B_sR_\infty^{\,T-1-s}\delta^T
+\text{higher order}.
\]

Assume $R_\infty$ is diagonalizable on the mean-zero subspace:
\[
R_\infty u_m = \mu_m u_m,
\qquad
|\mu_2|\ge\cdots\ge|\mu_n|,
\quad |\mu_m|<1.
\]
Expanding $\delta^T=\sum_{k=2}^n a_k\lambda_k^T v_k$ and $v_k=\sum_{m=2}^n c_{mk}u_m$ yields vector-level asymptotic
\[
h^0-g^0
=\sum_{k=2}^n\sum_{m=2}^n a_k c_{mk}\,(\lambda_k\mu_m)^T\,u_m
+\mathcal R_T,
\]
where $ \mathcal R_T$ is lower order under a spectral gap/non-resonance condition.
In particular, if $a_2c_{22}\neq 0$ and $|\lambda_2\mu_2|>\max_{(k,m)\neq(2,2)}|\lambda_k\mu_m|$, then
\[
h^0-g^0
= a_2c_{22}(\lambda_2\mu_2)^T u_2 + o((\lambda_2\mu_2)^T).
\]
Hence
\[
=\norm{h^0-g^0}
= \mathcal O\big((|\lambda_2|\,|\mu_2|)^T\big).
\]

\paragraph{Reversible symmetric case.}
If $P$ is reversible with respect to $h^\infty$, then asymptotically $R_\infty$ has the same nontrivial spectrum as $M$, so $|\mu_2|=|\lambda_2|$, yielding
\[
\norm{h^0-g^0}=\mathcal O\big(|\lambda_2|^{2T}\big).
\]
So initializing backward recursion at $h^\infty$ creates a second-order (in spectral decay) terminal error at time $0$.

\paragraph{Answer to ``depends only on $P$?''}
$R_\infty$ depends only on $P$ (through $h^\infty$), but the leading coefficient $a_2c_{22}$ depends on $h^0$.
So the decay \emph{rate} is determined by spectral quantities of $P$, while the leading amplitude of $h^0-g^0$ is not, in general, a function of $P$ alone.

\section{Practical interpretation}
\begin{itemize}
\item Forward warm-up error: $h^T-h^\infty$ decays like the first nontrivial forward mode.
\item Backward propagated error: $h^0-g^0$ equals terminal forward error filtered by the backward propagator.
\item Under stable reverse dynamics, this yields multiplicative decay rates and, in reversible settings, a squared spectral-decay effect.
\end{itemize}

\section{Designing $P$ for fastest decay of $h^0-g^0$}
Now assume $h^0$ is fixed and $P$ is designable (possibly under constraints). The previous sections show that the asymptotic rate is controlled by products $|\lambda_k\mu_m|$, and in reversible/similar settings by $|\lambda_2|^2$.

\subsection*{Unconstrained optimum}
If the only constraints are row-stochasticity and prescribed invariant law $h^\infty$,
the rank-one kernel
\[
P^\star=\ones (h^\infty)^\top
\]
is admissible and satisfies
\[
h^{t+1}=P^{\star\top}h^t=h^\infty,\qquad t\ge 0.
\]
Hence $h^T=h^\infty$ for $T\ge 1$, $Q^t=P^\star$, and therefore
\[
h^0-g^0=0\qquad (T\ge 1).
\]
So without structural constraints, the best possible asymptotic (indeed finite-time) rate is immediate convergence.

\subsection*{Constrained design principle}
In practice one often imposes structure (sparsity/locality, bounded jumps, reversibility, parametric form). Then exact one-step mixing is typically impossible.
The dominant asymptotic factor becomes
\[
\rho_{\mathrm{eff}}(P)\coloneqq \max_{k,m\ge 2} |\lambda_k(P)\,\mu_m(P)|,
\]
and (under mode non-degeneracy) 
\[
\norm{h^0-g^0}=\Theta\!\big(\rho_{\mathrm{eff}}(P)^T\big).
\]
Thus a natural design problem is
\[
\min_{P\in\mathcal C}\ \rho_{\mathrm{eff}}(P),
\]
where $\mathcal C$ is the admissible kernel class.

\subsection*{Reversible case}
If $P$ is constrained to be reversible w.r.t.\ $h^\infty$, then $R_\infty$ is similar to $P$ and $|\mu_2|=|\lambda_2|$, giving
\[
\norm{h^0-g^0}=\Theta\!\big(|\lambda_2(P)|^{2T}\big)
\]
for generic initial projection. Therefore fastest asymptotic convergence is obtained by maximizing spectral gap
\[
\operatorname{gap}(P)=1-|\lambda_2(P)|
\]
inside the constraint set:
\[
\max_{P\in\mathcal C}\ \operatorname{gap}(P)
\quad\Longleftrightarrow\quad
\min_{P\in\mathcal C} |\lambda_2(P)|.
\]

\subsection*{Dependence on $h^0$}
For fixed $P$, the prefactor depends on projection of $h^0-h^\infty$ on slow modes. If one can adapt $P$ to $h^0$ (e.g. finite-horizon design), one may additionally optimize this prefactor, but asymptotically the exponential rate is still governed by the spectral objective above.

\end{document}
